
% ============================================================
% VIABILITY PHILOSOPHY SYSTEM
% BOOK I: ONTOLOGY OF VIABILITY
%
% Project: THE_PHILOSOPY_OF_VIABILITY_LIBRARY_PROJECT
% File: VS.B01_ONTOLOGY_OF_VIABILITY/macros.tex
% Location: VS.B01_ONTOLOGY_OF_VIABILITY/
%
% Purpose:
% Central collection of project-specific LaTeX commands
% and environments for:
%
% - System and book identification
% - Structural identifiers
% - Logical entity identifiers
% - Internal cross-referencing
% - Formal entity presentation
% - Mathematical and symbolic notation
% - Ontological terminology
% - Argument and proof presentation
% - Editorial and explanatory notes
% - Supplementary material references
%
% Compiler: pdfLaTeX
%
% Dependencies:
% - amsmath
% - amssymb
% - hyperref
% - cleveref
%
% These packages are loaded in preamble.tex.
%
% No identifier is automatically assigned by this file.
% No logical, ontological, or epistemological commitment
% is established merely through the use of a macro.
%
% ============================================================


% ============================================================
% 1. GENERAL SYSTEM AND BOOK COMMANDS
% ============================================================

\newcommand{\ViabilitySystem}{Viability Philosophy System}

\newcommand{\ViabilityBook}{Ontology of Viability}

\newcommand{\VS}{VS}

\newcommand{\BOne}{B01}


% ============================================================
% 2. STRUCTURAL IDENTIFIER COMMANDS
%
% Structural identifiers specify the organizational location
% of a textual or structural unit within the complete system.
%
% Example:
% VS.B01.P01.C01.S01.SS01.P001
%
% Structural identifiers:
% - Are supplied explicitly.
% - Are not generated by these commands.
% - Do not represent logical entity identifiers.
% - Must conform to the project's structural registry.
%
% These commands only provide consistent presentation.
% ============================================================

\newcommand{\StructuralIdentifier}[1]{%
    \texttt{#1}%
}

\newcommand{\BookStructuralIDFormat}[1]{%
    \StructuralIdentifier{#1}%
}

\newcommand{\PartStructuralID}[1]{%
    \StructuralIdentifier{#1}%
}

\newcommand{\ChapterStructuralID}[1]{%
    \StructuralIdentifier{#1}%
}

\newcommand{\SectionStructuralID}[1]{%
    \StructuralIdentifier{#1}%
}

\newcommand{\SubsectionStructuralID}[1]{%
    \StructuralIdentifier{#1}%
}

\newcommand{\ParagraphStructuralID}[1]{%
    \StructuralIdentifier{#1}%
}


% ============================================================
% 3. LOGICAL ENTITY IDENTIFIER COMMANDS
%
% Logical entity identifiers specify the registered identity
% of formal entities within the Viability Philosophy System.
%
% Entity identifiers are intended to remain globally stable
% across books and structural locations.
%
% Approved prefixes in the current identifier convention:
%
% AX = Axiom
% DF = Definition
% PR = Proposition
% TH = Theorem
% HY = Hypothesis
% AR = Argument
% PF = Proof
%
% The numerical component is supplied explicitly.
%
% These commands do not:
% - Assign identifiers.
% - Verify registry membership.
% - Establish logical validity.
% - Establish proof or derivability.
% ============================================================

\newcommand{\LogicalEntityIdentifier}[1]{%
    \texttt{#1}%
}

\newcommand{\EntityID}[2]{%
    \LogicalEntityIdentifier{#1-#2}%
}

\newcommand{\AxiomID}[1]{%
    \EntityID{AX}{#1}%
}

\newcommand{\DefinitionID}[1]{%
    \EntityID{DF}{#1}%
}

\newcommand{\PropositionID}[1]{%
    \EntityID{PR}{#1}%
}

\newcommand{\TheoremID}[1]{%
    \EntityID{TH}{#1}%
}

\newcommand{\HypothesisID}[1]{%
    \EntityID{HY}{#1}%
}

\newcommand{\ArgumentID}[1]{%
    \EntityID{AR}{#1}%
}

\newcommand{\ProofID}[1]{%
    \EntityID{PF}{#1}%
}


% ============================================================
% 4. IDENTIFIED LOGICAL ENTITY PRESENTATION
%
% These commands display an entity's type, identifier,
% and descriptive content.
%
% They do not create theorem environments.
% They do not assign identifiers.
% They do not establish logical dependencies.
% They do not replace the corresponding registers.
%
% The descriptive content may contain ordinary LaTeX
% mathematical expressions and formatting commands.
% ============================================================

\newcommand{\IdentifiedAxiom}[2]{%
    \par\noindent
    \textbf{Axiom \AxiomID{#1}.}%
    \enspace #2%
    \par
}

\newcommand{\IdentifiedDefinition}[2]{%
    \par\noindent
    \textbf{Definition \DefinitionID{#1}.}%
    \enspace #2%
    \par
}

\newcommand{\IdentifiedProposition}[2]{%
    \par\noindent
    \textbf{Proposition \PropositionID{#1}.}%
    \enspace #2%
    \par
}

\newcommand{\IdentifiedTheorem}[2]{%
    \par\noindent
    \textbf{Theorem \TheoremID{#1}.}%
    \enspace #2%
    \par
}

\newcommand{\IdentifiedHypothesis}[2]{%
    \par\noindent
    \textbf{Hypothesis \HypothesisID{#1}.}%
    \enspace #2%
    \par
}

\newcommand{\IdentifiedArgument}[2]{%
    \par\noindent
    \textbf{Argument \ArgumentID{#1}.}%
    \enspace #2%
    \par
}


% ============================================================
% 5. INTERNAL STRUCTURAL CROSS-REFERENCES
%
% These commands refer to LaTeX labels.
%
% The referenced labels must be explicitly defined
% in the relevant document files.
%
% The commands do not generate labels automatically.
% ============================================================

\newcommand{\StructuralRef}[1]{%
    \cref{#1}%
}

\newcommand{\StructuralRefUpper}[1]{%
    \Cref{#1}%
}

\newcommand{\PageRef}[1]{%
    page~\pageref{#1}%
}

\newcommand{\ChapterRef}[1]{%
    Chapter~\ref{#1}%
}

\newcommand{\PartRef}[1]{%
    Part~\ref{#1}%
}

\newcommand{\SectionRef}[1]{%
    Section~\ref{#1}%
}

\newcommand{\SubsectionRef}[1]{%
    Subsection~\ref{#1}%
}

\newcommand{\EquationRef}[1]{%
    \eqref{#1}%
}

\newcommand{\FigureRef}[1]{%
    \cref{#1}%
}

\newcommand{\TableRef}[1]{%
    \cref{#1}%
}

\newcommand{\TheoremRef}[1]{%
    \cref{#1}%
}

\newcommand{\DefinitionRef}[1]{%
    \cref{#1}%
}

\newcommand{\PropositionRef}[1]{%
    \cref{#1}%
}


% ============================================================
% 6. LOGICAL ENTITY REFERENCES
%
% These commands display the global identifier of
% a logical entity.
%
% An entity reference is not equivalent to:
% - A structural location reference.
% - A proof of the entity.
% - A claim that the entity is established.
% - A claim that the entity has a valid registry entry.
%
% Registry membership and logical status must be verified
% independently.
% ============================================================

\newcommand{\AxiomRef}[1]{%
    \AxiomID{#1}%
}

\newcommand{\DefinitionEntityRef}[1]{%
    \DefinitionID{#1}%
}

\newcommand{\PropositionEntityRef}[1]{%
    \PropositionID{#1}%
}

\newcommand{\TheoremEntityRef}[1]{%
    \TheoremID{#1}%
}

\newcommand{\HypothesisRef}[1]{%
    \HypothesisID{#1}%
}

\newcommand{\ArgumentRef}[1]{%
    \ArgumentID{#1}%
}

\newcommand{\ProofRef}[1]{%
    \ProofID{#1}%
}


% ============================================================
% 7. FORMAL SYMBOL AND EXPRESSION PRESENTATION
%
% General-purpose commands for mathematical and formal
% notation.
%
% These commands provide mathematical presentation only.
% Their semantic interpretation must be established
% in the relevant formal or philosophical context.
% ============================================================

\newcommand{\FormalSymbol}[1]{%
    \ensuremath{#1}%
}

\newcommand{\FormalExpression}[1]{%
    \ensuremath{#1}%
}

\newcommand{\SymbolicDefinition}[2]{%
    \ensuremath{#1\;:=\;#2}%
}


% ============================================================
% 8. TERMINOLOGY AND CONCEPTUAL PRESENTATION
% ============================================================

\newcommand{\TechnicalTerm}[1]{%
    \textit{#1}%
}

\newcommand{\Concept}[1]{%
    \textit{#1}%
}

\newcommand{\DefinedTerm}[1]{%
    \textbf{#1}%
}

\newcommand{\EmphasizedTerm}[1]{%
    \emph{#1}%
}

% Displays a term and adds it to the index when the
% \index command is available.
%
% This conditional form allows the macro file to remain
% usable even when the indexing package is not loaded.
%
% Index generation itself requires an appropriate
% indexing configuration and compilation workflow.

\newcommand{\TermIndex}[1]{%
    #1%
    \ifcsname index\endcsname
        \index{#1}%
    \fi
}


% ============================================================
% 9. PRINCIPAL ONTOLOGICAL TERMINOLOGY
%
% These commands provide consistent typographical
% presentation of principal terms used in Book I.
%
% They do not:
% - Define the philosophical meaning of these terms.
% - Establish relations among their referents.
% - Introduce additional ontological commitments.
%
% Their conceptual definitions belong to the relevant
% sections and formal registers of the book.
% ============================================================

\newcommand{\Identity}{%
    \textit{Identity}%
}

\newcommand{\Manifestation}{%
    \textit{Manifestation}%
}

\newcommand{\Being}{%
    \textit{Being}%
}

\newcommand{\RealmsOfBeing}{%
    \textit{Realms of Being}%
}

\newcommand{\OntologicalLight}{%
    \textit{Ontological Light}%
}

\newcommand{\DynamicsOfBeing}{%
    \textit{Dynamics of Being}%
}

\newcommand{\FirstDetermination}{%
    \textit{First Determination}%
}

\newcommand{\StructuralDetermination}{%
    \textit{Structural Determination}%
}

\newcommand{\RelationalManifestation}{%
    \textit{Relational Manifestation}%
}

\newcommand{\MaterialManifestation}{%
    \textit{Material Manifestation}%
}

\newcommand{\HubOfDecree}{%
    \textit{Hub of Decree}%
}

\newcommand{\HubOfSignification}{%
    \textit{Hub of Signification}%
}

\newcommand{\HubOfRealization}{%
    \textit{Hub of Realization}%
}


% ============================================================
% 10. ARGUMENT AND PROOF PRESENTATION
% ============================================================

\newcommand{\ArgumentLabel}[1]{%
    \par\noindent
    \textbf{Argument:}%
    \enspace #1%
    \par
}

\newcommand{\ProofLabel}[1]{%
    \par\noindent
    \textbf{Proof:}%
    \enspace #1%
    \par
}

\newcommand{\PremiseLabel}[1]{%
    \textbf{Premise:}%
    \enspace #1%
}

\newcommand{\ConclusionLabel}[1]{%
    \textbf{Conclusion:}%
    \enspace #1%
}


% ============================================================
% 11. RELATIONAL AND DEPENDENCY NOTATION
%
% General-purpose symbolic notation.
%
% The semantic meaning of each relation must be explicitly
% established in its relevant philosophical or formal context.
%
% In particular, the symbols do not independently establish
% ontological priority, causal dependence, or temporal order.
% ============================================================

\newcommand{\DependsOn}{%
    \ensuremath{\longrightarrow}%
}

\newcommand{\Entails}{%
    \ensuremath{\vdash}%
}

\newcommand{\Implies}{%
    \ensuremath{\Rightarrow}%
}

\newcommand{\EquivalentTo}{%
    \ensuremath{\Leftrightarrow}%
}

\newcommand{\NotEquivalentTo}{%
    \ensuremath{\not\Leftrightarrow}%
}

\newcommand{\Necessary}{%
    \ensuremath{\Box}%
}

\newcommand{\Possible}{%
    \ensuremath{\Diamond}%
}


% ============================================================
% 12. EDITORIAL AND EXPLANATORY ENVIRONMENTS
%
% These environments provide typographical containers
% for explanatory material.
%
% They do not establish the epistemic or ontological
% status of the content they contain.
% ============================================================

\newenvironment{editorialnote}
{%
    \begin{quote}
    \small
    \noindent
    \textbf{Editorial Note.}\enspace
}
{%
    \end{quote}
}

\newenvironment{technicalnote}
{%
    \begin{quote}
    \small
    \noindent
    \textbf{Technical Note.}\enspace
}
{%
    \end{quote}
}

\newenvironment{conceptualnote}
{%
    \begin{quote}
    \small
    \noindent
    \textbf{Conceptual Note.}\enspace
}
{%
    \end{quote}
}


% ============================================================
% 13. QUOTATIONS AND SOURCE DISCUSSION
% ============================================================

\newenvironment{sourcequotation}
{%
    \begin{quote}
    \small
}
{%
    \end{quote}
}

\newcommand{\SourceDiscussion}[1]{%
    \par\noindent
    \textbf{Source Discussion.}%
    \enspace #1%
    \par
}


% ============================================================
% 14. APPENDIX AND SUPPLEMENTARY MATERIAL REFERENCES
% ============================================================

\newcommand{\AppendixRef}[1]{%
    Appendix~\ref{#1}%
}

\newcommand{\SupplementaryMaterialRef}[1]{%
    \cref{#1}%
}


% ============================================================
% 15. CONSISTENT LABEL PREFIX CONVENTIONS
%
% Recommended label naming conventions:
%
% ch:       Chapters
% part:     Parts
% sec:      Sections
% subsec:   Subsections
% eq:       Equations
% fig:      Figures
% tab:      Tables
% ax:       Axioms
% df:       Definitions
% pr:       Propositions
% th:       Theorems
% hy:       Hypotheses
% ar:       Arguments
% pf:       Proofs
% app:      Appendices
%
% These prefixes are organizational conventions only.
%
% They do not automatically create labels or guarantee
% uniqueness. Every label must be explicitly defined.
% ============================================================


% ============================================================
% END OF macros.tex
% ============================================================
